
Par
ailleurs, on a $K=\varinjlim_f A_f$, où $f$ parcourt l'ensemble des éléments non nuls de $A$. Par suite, $\Spec(K)=\varprojlim_f\Spec(A_f)$, et :
\[
X_s=X\times_S\Spec(K)=\varprojlim_f X\times_S\Spec(A_f)
\]
d'après EGA IV 8.2.3 et 8.2.5. Comme les schémas envisagés sont noethériens, et les morphismes de transition affines et plats, on a d'après IV 3.2.3 :
\[
K_\bullet(X_s)=\varinjlim_f K_\bullet(X\times_S\Spec(A_f)).
\]
Comme les homomorphismes «~image inverse~» sur les $K_\bullet$ sont linéaires par rapport aux $K^{\bullet}$ (IV (2.12.2)), il existe un ouvert non vide $U=\Spec(A_f)$ de $S$ tel que dans $K_\bullet(X_U)$ on ait :
\[
y_U=\sum_i\xi_i\cdot n_i.
\]
Soit $S'$ un sous-schéma fermé de $S$ de support $S-U$. Alors $\varphi_{S'}$ est surjective d'après l'hypothèse de récurrence, et par suite le diagramme (1.2.1) montre que :
\[
y-\sum_i\xi_i\cdot n_i\in \operatorname{Im}(\varphi),
\]
donc, comme $n_i=f^*(n_i)$, $\varphi$ est surjective.

\textbf{Corollaire 1.3.} \uline{Soit} $f:X\to S$ \uline{un morphisme plat de schémas noethériens}. \uline{On suppose que pour tout point} $s$ \uline{de} $S$, \uline{l'homomorphisme} :
\[
K^{\bullet}(X)\longrightarrow K_\bullet(X_s)
\]
\uline{est surjectif}. \uline{Alors l'homomorphisme déduit de} $f^*$ :
\[
K_\bullet(S)\otimes_{K^{\bullet}(S)}K^{\bullet}(X)\longrightarrow K_\bullet(X)
\]
\uline{est surjectif}.

